% =========================================================
% Proof: Injectivity of Zermelo successor
% Source: volume-i/axiom-systems/notes/zermelo/zermelo-numerals.tex
% =========================================================

\subsection*{Injectivity of Zermelo Successor}
\label{prf:zermelo-successor-injective}

\begin{remark*}[Return]
\hyperref[prop:zermelo-successor-injective]{$\leftarrow$ Back to Proposition (Injectivity of Zermelo successor) in Notes}
\end{remark*}

\begin{proof}
\Claim For all sets $a, b$, $S_Z(a) = S_Z(b) \Rightarrow a = b$.

\Given The Zermelo successor $S_Z(n) = \{n\}$
(\hyperref[def:zermelo-successor]{Definition~(Zermelo successor)}). The
Proposition that singletons are well-defined
(\hyperref[prop:singleton-well-defined]{Singleton is well-defined}): the unique
element of $\{a\}$ is $a$.

\Goal To show: if $\{a\} = \{b\}$ then $a = b$.

\Strategy Suppose $\{a\} = \{b\}$. We extract the unique element of each side
and use the fact that equal sets have the same elements.

Since $\{a\} = \{b\}$, the sets have the same elements. We know $a \in \{a\}$
(by definition of singleton). Therefore $a \in \{b\}$. By definition of $\{b\}$,
the only element of $\{b\}$ is $b$. Hence $a = b$. \AsReq
\end{proof}

\begin{remark*}[Proof shape]
The argument is a one-step membership extraction. If two singletons are equal as
sets, then any element of one is an element of the other. Since $a$ is the
unique element of $\{a\}$ and also an element of $\{b\}$, and $b$ is the unique
element of $\{b\}$, we conclude $a = b$. The proof does not require the Axiom
of Foundation; it uses only the definition of singleton and set equality.
\end{remark*}

\begin{remark*}[Contrast with von Neumann successor]
The von Neumann successor $S_V(n) = n \cup \{n\}$ also satisfies injectivity,
but the proof there is more involved because $S_V(a)$ has more than one element
when $a \neq \varnothing$. For Zermelo numerals, injectivity is immediate
because each successor is a singleton.
\end{remark*}

\begin{remark*}[Dependencies]
The proof depends on: the Definition of Zermelo successor
(\hyperref[def:zermelo-successor]{Definition~(Zermelo successor)}) and
\hyperref[prop:singleton-well-defined]{Singleton is well-defined}.
\end{remark*}








% ============================================================
% PROOF RECORD
% ============================================================
% Proof Name: Upper Bound is Unique
% Dependency Claims:
%   - Definition of Upper Bound
% Difficulty: ★
% ============================================================

\begin{theorem}[Upper Bound is Unique]
Let $A \subseteq \mathbb{R}$.
If $u$ and $v$ are both upper bounds of $A$, then $u = v$.
\end{theorem}

\begin{proof}
\Given $u$ and $v$ are upper bounds of $A$.
\Goal To show that $u = v$.
\Strategy We proceed directly by unpacking the definition of upper bound.

\ByDef Since $u$ is an upper bound of $A$, we have
\[
\forall a \in A, \quad a \le u.
\] \tagDU

\ByDef Since $v$ is an upper bound of $A$, we have
\[
\forall a \in A, \quad a \le v.
\] \tagDU

% ------------------------------------------------------------
% Main argument goes here.
% You must use only:
%   - Definition of upper bound
%   - Order properties of ℝ
%   - Antisymmetry of ≤
% ------------------------------------------------------------

\Hence $u = v$. \AsReq
\end{proof}